% specimen-tikz.tex -- florilegium's species drawn in TikZ (\usetikzlibrary{florilegium}).
% Under pdfLaTeX pgfmath computes and TikZ draws (pure TikZ, the proof, about 6 minutes);
% under LuaLaTeX florilegium-organs.lua computes the same pgf commands and TikZ draws them
% (engine=lua, the library's default there): the same PDF, byte for byte, in about 20
% seconds. \pgfkeys{/florilegium/engine=tikz} draws in pure TikZ under LuaLaTeX too.
% Page 1: two margins' worth of plants, three species each, grown as florilegium-grow.lua
% grows a margin; page 2: each body plan's flower and leaf close up, and one plant through
% the seasons.
% Copyright (c) 2026 Kevin Raymond. MIT License.
%   lualatex specimen-tikz.tex
%   extra_mem_top=10000000 pdflatex specimen-tikz.tex
% pdfLaTeX keeps a page's drawing in main memory until the page is shipped out, and one
% strip of plants in v1.3's detail takes about 7 million words of it, more than pdfTeX's
% default 5 million: extra_mem_top raises it for the run. LuaLaTeX's memory grows as needed.
\documentclass[11pt]{article}
\usepackage[a4paper,margin=12mm]{geometry}
\usepackage{tikz}
\usetikzlibrary{florilegium}
\pagestyle{empty}
\setlength\parindent{0pt}
\newcommand\traits[2]{% #1 seed, #2 role: a few of the species' traits, for a caption
  \florilegiumspecies{#1}{#2}%
  \florilegiumtrait{plan} #1: \florilegiumtrait{merosity}~parts, \florilegiumtrait{form},
  \florilegiumtrait{inflo}, \florilegiumtrait{color}, for \florilegiumtrait{syndrome};
  \florilegiumtrait{compound} \florilegiumtrait{leafmargin} leaves, \florilegiumtrait{arrangement}}
\newcommand\closeup[2]{% #1 plan, #2 seed: its flower and leaf, and a caption
  \begin{tikzpicture}
    \pic at (1.4,1.6) {florilegium flower={seed=#2, role=#1, size=2.6cm}};
    \pic at (3.6,0) {florilegium leaf={seed=#2, role=#1, size=3.4cm}};
  \end{tikzpicture}\par
  \parbox{5.4cm}{\raggedright\scriptsize\traits{#2}{#1}.}}
\begin{document}
{\Large Species in TikZ}\par\smallskip
{\small Each plant is a species sampled from a genome (body plan, the visitor its flowers are
for, the flower, its cluster, the leaf, the fruit) and drawn by TikZ in florilegium's v1.3
hand. Under pdfLaTeX, pgfmath computes every point; under LuaLaTeX, Lua computes the same pgf
commands, to the byte, and sooner. Three species share a root, as a margin's main stems do: stems and
the flowers at their tops first, then leaves at the nodes, then branches wherever there is
room, a pale flower of the first behind them.}\par\medskip
\begin{tikzpicture}
  \pic at (0,0) {florilegium plants={plants={3/rose,17/umbellifer,29/mustard}, size=21cm, width=6cm}};
  \pic at (9.3,0) {florilegium plants={plants={3/daisy,17/lily,29/rose}, size=21cm, width=6cm}};
\end{tikzpicture}\par\medskip
{\scriptsize\begin{minipage}[t]{0.48\linewidth}\raggedright
  \traits{3}{rose}.\\ \traits{17}{umbellifer}.\\ \traits{29}{mustard}.\end{minipage}\hfill
\begin{minipage}[t]{0.48\linewidth}\raggedright
  \traits{3}{daisy}.\\ \traits{17}{lily}.\\ \traits{29}{rose}.\end{minipage}}
\newpage
{\large Each body plan close up}\par\medskip
\begin{minipage}[t]{0.32\linewidth}\closeup{rose}{17}\end{minipage}\hfill
\begin{minipage}[t]{0.32\linewidth}\closeup{mustard}{3}\end{minipage}\hfill
\begin{minipage}[t]{0.32\linewidth}\closeup{daisy}{29}\end{minipage}\par\bigskip
\begin{minipage}[t]{0.32\linewidth}\closeup{umbellifer}{3}\end{minipage}\hfill
\begin{minipage}[t]{0.32\linewidth}\closeup{lily}{29}\end{minipage}\hfill
\begin{minipage}[t]{0.32\linewidth}\closeup{rose}{29}\end{minipage}\par\bigskip
{\large One plant through the year}\par\medskip
\begin{tikzpicture}
  \foreach \a/\name [count=\k from 0] in {0.3/spring,0.75/summer,1.25/autumn}{
    \pic at (\k*6.2,0) {florilegium plant={seed=17, role=rose, size=8.5cm, width=5.2cm, season=\a}};
    \node[font=\footnotesize] at (\k*6.2,-0.4) {\name};
  }
\end{tikzpicture}\par
{\footnotesize Rose 17 in spring, summer and autumn: its flowers open from the top of the
plant down, so the top is the first to open and the first to turn to hips.}
\end{document}
